\section{Background}\label{sec:background}

\subsection{\texorpdfstring{$\LTLf$}{LTLf}}
\guide{Syntax and finite-trace semantics; strong/weak next ($X[!]$ vs.\ $X$); notation for
traces over $2^{\mathcal{I}\cup\mathcal{O}}$. Cite \cite{DeGiacomoVardi2013}.}

$\LTLf$~\cite{DeGiacomoVardi2013} is a variant of LTL that is interpreted over finite traces.
That means we have the same operators as LTL, such as $G$ (Globally), $F$ (Finally), $X$ (Next), etcetera.
Traces must be of length at least 1.
A trace is a word over the alphabet $2^{\mathcal{V}}$, where $\mathcal{V}$ is the set of variables; we write $\pi[i]$ for its letter at step $i$ and $\pi[0..i]$ for its prefix up to step $i$.
For a letter $v$ and a set of variables $U$, $v|_U = v \cap U$ is the restriction of $v$ to $U$.

The Next operator usually assumes that there is always a next state; on finite traces there is not, so we have weak and strong next operators.
The weak next operator $X$ is true in the last state of a trace, while the strong next operator $X[!]$ is false in the last state of a trace.

A similar logic is Past-$\LTLf$, which is interpreted over finite traces but with past operators.
Each future operator has a past counterpart: $Y$ (Yesterday) for $X$, $O$ (Once) for $F$, $H$ (Historically) for $G$, and $S$ (Since) for $U$ (Until).

\subsection{Automata on Finite Words}
\guide{NFA, DFA, determinisation blow-up, minimisation. Past-$\LTLf$ and the reverse/mirror
construction that gives a small NFA (needed for Method 1).}

$\LTLf$ formulas can be translated into automata on finite words that accept exactly the traces that satisfy the formula.
The automata can be non-deterministic (NFA) or deterministic (DFA).
The NFA of a formula can have $2^{O(|\varphi|)}$ states, and the DFA $2^{2^{O(|\varphi|)}}$ states.

The translation from $\LTLf$ to DFA can be done by translating to NFA and determinizing, or directly by \emph{formula progression}: each state is a formula, and reading a letter rewrites it into the formula that the rest of the trace must satisfy.

Using Past-$\LTLf$, we can get an NFA of single-exponential size another way.
We replace every operator of $\varphi$ with its past counterpart, translate the resulting past formula to a DFA, which is only single-exponential for past formulas, and then reverse that DFA to get an NFA for $\varphi$.

\subsection{Reactive Synthesis on Finite Traces}
\guide{Controller as a transducer; realizability; DFA game.
Turn order (who moves first). Cite \cite{PnueliRosner1989,DeGiacomoVardi2015}.}

$\LTLf$ synthesis~\cite{PnueliRosner1989,DeGiacomoVardi2015} is a problem in which we are given a specification $\varphi$ and sets of input and output variables $\mathcal{I}$ and $\mathcal{O}$.
We interpret the specification as a game between the system and the environment, in which the environment chooses values for the input variables, and the system chooses values for the output variables.
On finite traces the system also decides when the trace ends: it may stop after any step.
The system wins if the trace it stops on satisfies the specification, and the synthesis problem is realizable if there exists a strategy for the system to win against every possible strategy of the environment.

Such a strategy can be implemented as a transducer, a finite state machine that reads the input variables, stores the history of the trace in its state, and outputs values for the output variables.

A nuance of the synthesis problem is the turn order, which determines who moves first.
If the environment moves first in each step, the system sees the current inputs before choosing its outputs; this is a \emph{Mealy} controller.
If the system moves first, it sees only past inputs; this is a \emph{Moore} controller.
We use the Mealy setting throughout.

\subsection{Symbolic Representations}
\guide{BDDs, multi-terminal DFAs (MTDFA) in Spot; why an mtdfa representation is a second
implementation of a method, not a new method. Cite \cite{Spot2022}.}

Storing the automata as a list of states and transitions is called an explicit representation.
It is intuitive but each letter of the alphabet gets its own edge, and the alphabet is exponential in the number of variables.
Symbolic representations instead group letters into Boolean functions.

A binary decision diagram (BDD) represents a Boolean function over the variables, and so a set of valuations: those on which the function is true.
A multi-terminal BDD (MTBDD) generalizes this to functions whose result is one of several terminal values instead of true or false.
Spot~\cite{Spot2022} uses MTBDDs to store a DFA as an \emph{MTDFA}: each state has one MTBDD whose terminals are the successor states, so all the transitions of a state are stored at once.

\subsection{Baseline Tool: \tool{ltlfsynt}}
\guide{What Spot's \tool{ltlfsynt} does, and how it is used here: as a baseline and as an
external realizability oracle.}

A baseline tool for $\LTLf$ synthesis is Spot's \tool{ltlfsynt}, which builds the MTDFA of the formula on the fly and solves the game as it goes.
We use it to compare the performance of our methods, and to check their verdicts.
